\documentclass[11pt,letterpaper]{article}
\usepackage[margin=0.75in]{geometry}
\usepackage[utf8]{inputenc}
\usepackage[T1]{fontenc}
\usepackage{lmodern}
\usepackage{helvet}
\usepackage{textcomp}
\usepackage{parskip}
\usepackage{longtable}
\usepackage{booktabs}
\usepackage{array}
\usepackage{calc}
\usepackage{amssymb}
\usepackage{enumitem}
\usepackage{titlesec}
\usepackage{xcolor}
\usepackage[colorlinks=true,urlcolor=blue,linkcolor=black]{hyperref}
\renewcommand{\familydefault}{\sfdefault}
\setlist[itemize]{leftmargin=1.3em,itemsep=1pt,topsep=2pt}
\setlist[enumerate]{leftmargin=1.5em,itemsep=1pt,topsep=2pt}
\titleformat{\section}{\large\bfseries}{}{0em}{}[\vspace{0.1em}\hrule\vspace{0.2em}]
\newcolumntype{L}[1]{>{\raggedright\arraybackslash}p{#1}}
\providecommand{\tightlist}{\setlength{\itemsep}{0pt}\setlength{\parskip}{0pt}}
\setlength{\parindent}{0pt}\setlength{\parskip}{4pt}
\begin{document}
\section*{Teaching Statement --- \'{A}ngel Luis Robles-Fern\'andez}
\textbf{Oregon State University --- Assistant Professor, Integrative Biology}

My goal is to help students connect biological observations, theory, and quantitative evidence so they can formulate and answer questions independently. My training in physics and evolutionary biology, together with experience in industrial machine learning, informs a teaching approach in which computation serves biological understanding rather than replacing organismal observation, experiments, or natural history.

\textbf{Experience and teaching practice.}
I have taught Introduction to R, served as a teaching assistant for Quantitative Ecology and Statistical Models for Biology at Arizona State University, and been invited to teach graduate-level ecological niche modeling at INECOL. As an RStudio Certified Trainer in the Tidyverse, I use live coding, worked examples, and collaborative debugging. In niche-modeling instruction, I scaffolded projects from occurrence-data cleaning through model fitting and interpretation. Students practiced a complete scientific workflow rather than learning isolated software commands.

\textbf{Inclusive, evidence-based learning.}
I will combine short explanations with active learning, peer discussion, and low-stakes exercises that reveal misconceptions before major assessments. Early, ungraded diagnostics will help me adapt support to differences in biological and programming preparation. Annotated examples and progressively less guided tasks will give beginners an entry point while allowing experienced students to extend their analyses. Small-group activities will rotate roles so that coding experience does not determine who participates. I will use open resources and modest-sized datasets, avoiding requirements for paid AI subscriptions. Feedback from students and formative assessments will guide revisions to pacing and instruction.

\textbf{Biological reasoning and AI literacy.}
Assignments will begin with a biological question and require students to justify data, assumptions, and interpretation. An advanced project can use small or simulated datasets to compare Bayesian XSDM models with reference priors and experimentally informed priors. Students will explain how performance evidence changes predictions and uncertainty, and what assumptions connect an experimental response to population persistence. My open-source tools will support learning, not replace biological reasoning. Where AI assistance is permitted, students will disclose its use, verify sources, test generated code, and explain their choices. Reproducible notebooks and project checkpoints will assess understanding rather than polished output alone.

\textbf{Mentoring and departmental contributions.}
I will establish clear expectations, regular meetings, and transparent contribution-based credit for undergraduate and graduate researchers. New students will begin with supported data and analysis tasks, then take ownership of a question; graduate students will develop independent methods and interdisciplinary collaborations. Mentoring plans will be adapted to students' preparation, circumstances, and career goals. My industry experience will help me discuss careers beyond academia alongside preparation for graduate research.

I am prepared to contribute to introductory biology, ecology, evolution, and quantitative instruction within the biology and zoology programs, and to offer Global Change Ecology or computational biology electives. In every setting, my objective is the same: students should leave able to connect a biological question to evidence, evaluate uncertainty, and communicate a defensible conclusion.
\end{document}
